\subsection{Iterated Tensor Products and Tensor Spaces}

The tensor-product universal property extends from bilinear maps to multilinear maps.\\
Its iterated form leads naturally to tensor spaces and contractions.

\begin{definition} \label{def:multilinear-map} \hlt{Multilinear Maps}\\
Let $n\in\mathbb{Z}_{>0}$ and let $\mathcal{V}_1,\ldots,\mathcal{V}_n,\mathcal{W}$ be $\mathbb{F}$-vector spaces.
\begin{enumerate}[label=(\roman*)]
\setlength{\itemsep}{0pt}
\item Fix $k\in\{1,\ldots,n\}$ and all entries except the $k$th. For $f:\mathcal{V}_1\times\cdots\times\mathcal{V}_n\to\mathcal{W}$, set $f_k(\mathbf{v})\defeq f(\mathbf{v}_1,\ldots,\mathbf{v},\ldots,\mathbf{v}_n)$.\\
The map $f$ is \hlt{multilinear}, or \hlt{$n$-linear}, if $f_k\in\LL(\mathcal{V}_k,\mathcal{W})$ in this setup.
\item The set of such maps is \hlt{$\operatorname{Mult}_{\mathbb{F}}(\mathcal{V}_1,\ldots,\mathcal{V}_n;\mathcal{W})$}.
\item If $\mathcal{W}=\mathbb{F}$, these maps are \hlt{multilinear forms}.
\end{enumerate}
\end{definition}

\begin{example} \label{ex:basic-multilinear-maps} \hlt{Basic Multilinear Maps}
\begin{enumerate}[label=(\roman*)]
\setlength{\itemsep}{0pt}
\item Let $A$ be an associative $\mathbb{F}$-algebra and let $n\in\mathbb{Z}_{>0}$. Set $m_n(a_1,\ldots,a_n)\defeq a_1\cdots a_n$.\\
Distributivity and homogeneity in $A$ prove that $m_n:A^n\to A$ is $n$-linear.
\item Suppose $\dim\mathcal{V}=n\geq1$ and $\beta$ is an ordered basis.\\
The determinant of the $\beta$-coordinate columns is $n$-linear by Theorem \ref{thm:determinant-basic-laws}(ii). It is alternating by parts (iii)--(iv).
\end{enumerate}
\end{example}

\begin{definition} \label{def:iterated-tensor-product-universal-property} \hlt{Iterated Tensor Product by Universal Property}\\
Let $n\in\mathbb{Z}_{>0}$ and let $\mathcal{V}_1,\ldots,\mathcal{V}_n$ be $\mathbb{F}$-vector spaces.
\begin{enumerate}[label=(\roman*)]
\setlength{\itemsep}{0pt}
\item A pair $(\mathcal{T},t)$ is an \hlt{iterated tensor product} of $\mathcal{V}_1,\ldots,\mathcal{V}_n$ if both conditions hold.
\begin{enumerate}[label=(\alph*)]
\setlength{\itemsep}{0pt}
\item $t\in\operatorname{Mult}_{\mathbb{F}}(\mathcal{V}_1,\ldots,\mathcal{V}_n;\mathcal{T})$.
\item Let $\mathcal{Y}$ be an $\mathbb{F}$-vector space. Set $\mathcal{M}_{\mathcal{Y}}\defeq\operatorname{Mult}_{\mathbb{F}}(\mathcal{V}_1,\ldots,\mathcal{V}_n;\mathcal{Y})$. $(\forall\,f\in\mathcal{M}_{\mathcal{Y}})(\exists!\,\widetilde{f}\in\LL(\mathcal{T},\mathcal{Y}))\ f=\widetilde{f}\circ t$.
\end{enumerate}
\item Write $\mathcal{T}=\mathcal{V}_1\otimes\cdots\otimes\mathcal{V}_n$. The notation $\mathbf{v}_1\otimes\cdots\otimes\mathbf{v}_n$ abbreviates $t(\mathbf{v}_1,\ldots,\mathbf{v}_n)$.
\end{enumerate}
\end{definition}

\begin{illustrationfigure}[H]
  \centering
  % Source note: Roman, Advanced Linear Algebra, Ch. 14, Fig. 14.8.
% Redrawn as the universal-property diagram for iterated tensor products.
\begin{tikzpicture}[
  >=Stealth,
  every node/.style={font=\small},
  arrow/.style={->, line width=0.55pt},
  factor/.style={->, dashed, line width=0.55pt},
  note/.style={font=\scriptsize, text=gray!70!black}
]
  \node (prod) at (0,1.25) {$\mathcal{V}_1\times\cdots\times\mathcal{V}_n$};
  \node (tensor) at (4.25,1.25) {$\mathcal{V}_1\otimes\cdots\otimes\mathcal{V}_n$};
  \node (w) at (4.25,-0.8) {$\mathcal{W}$};

  \draw[arrow] (prod) -- node[above, note] {$t$ multilinear} (tensor);
  \draw[arrow] (prod) -- node[below left, note] {$f$ multilinear} (w);
  \draw[factor] (tensor) -- node[right, note] {$\widetilde{f}$ linear} (w);
\end{tikzpicture}

  \caption{The iterated tensor product is universal for multilinear maps.}
  \label{fig:multilinear-tensor-universal-property}
\end{illustrationfigure}

Figure \ref{fig:multilinear-tensor-universal-property} depicts the factorization in Definition \ref{def:iterated-tensor-product-universal-property}.\\
The definition, rather than the diagram, owns the universal property.

\begin{theorem} \label{thm:iterated-tensor-products-basic-laws} \hlt{Construction and Basic Laws for Iterated Tensor Products}\\
Let $n\in\mathbb{Z}_{>0}$ and let $\mathcal{V}_1,\ldots,\mathcal{V}_n,\mathcal{V},\mathcal{W}$ be $\mathbb{F}$-vector spaces. For $1\leq a\leq b\leq n$, set $\mathcal{V}_{a:b}^{\otimes}\defeq\mathcal{V}_a\otimes\cdots\otimes\mathcal{V}_b$.\\
Write $\mathbf{v}_{a:b}^{\otimes}\defeq\mathbf{v}_a\otimes\cdots\otimes\mathbf{v}_b$. Set $\mathcal{M}_{\mathcal{W}}\defeq\operatorname{Mult}_{\mathbb{F}}(\mathcal{V}_1,\ldots,\mathcal{V}_n;\mathcal{W})$.
\begin{enumerate}[label=(\roman*)]
\setlength{\itemsep}{0pt}
\item The following construction and uniqueness statements hold.
\begin{enumerate}[label=(\alph*)]
\setlength{\itemsep}{0pt}
\item Every binary parenthesization of $\mathcal{V}_1\otimes\cdots\otimes\mathcal{V}_n$ is an iterated tensor product.
\item Every iterated tensor product is spanned by its decomposable tensors.
\item Any two iterated tensor products are uniquely isomorphic by an isomorphism preserving decomposable tensors.
\end{enumerate}
\item Pointwise operations make $\mathcal{M}_{\mathcal{W}}$ an $\mathbb{F}$-vector space. Map $\Phi:\mathcal{M}_{\mathcal{W}}\to\LL(\mathcal{V}_{1:n}^{\otimes},\mathcal{W})$ defined by $\Phi(f)\defeq\widetilde{f}$ and is an isomorphism.
\item If $1\leq r<n$, set $\mathcal{A}_r\defeq\mathcal{V}_{1:r}^{\otimes}\otimes\mathcal{V}_{r+1:n}^{\otimes}$. The canonical isomorphism $\alpha_r:\mathcal{A}_r\xrightarrow{\sim}\mathcal{V}_{1:n}^{\otimes}$ satisfies $\alpha_r(\mathbf{v}_{1:r}^{\otimes}\otimes\mathbf{v}_{r+1:n}^{\otimes})=\mathbf{v}_{1:n}^{\otimes}$.
\item If $\pi\in S_n$, set $\mathcal{C}_\pi\defeq\mathcal{V}_{\pi(1)}\otimes\cdots\otimes\mathcal{V}_{\pi(n)}$. The canonical isomorphism $\tau_\pi:\mathcal{V}_{1:n}^{\otimes}\xrightarrow{\sim}\mathcal{C}_\pi$ satisfies $\tau_\pi(\mathbf{v}_{1:n}^{\otimes})=\mathbf{v}_{\pi(1)}\otimes\cdots\otimes\mathbf{v}_{\pi(n)}$.
\item The isomorphism $\lambda:\mathbb{F}\otimes\mathcal{V}\xrightarrow{\sim}\mathcal{V}$ is determined by $\lambda(r\otimes\mathbf{v})=r\mathbf{v}$. The isomorphism $\rho:\mathcal{V}\otimes\mathbb{F}\xrightarrow{\sim}\mathcal{V}$ is determined by $\rho(\mathbf{v}\otimes r)=r\mathbf{v}$.
\item If $(\forall\,i\in\{1,\ldots,n\})\,\dim\mathcal{V}_i<\infty$, then $\dim(\mathcal{V}_{1:n}^{\otimes})=\prod\limits_{i=1}^{n}\dim(\mathcal{V}_i)$.
\end{enumerate}
\begin{proof}
\begin{enumerate}[label=(\roman*)]
\setlength{\itemsep}{0pt}
\item
\begin{enumerate}[label=(\alph*)]
\setlength{\itemsep}{0pt}
\item For $n=1$, take $\mathcal{T}=\mathcal{V}_1$ and $t=I_{\mathcal{V}_1}$. The unique mediator of a linear map $f:\mathcal{V}_1\to\mathcal{Y}$ is $f$ itself.\\
Assume the assertion for fewer than $n$ factors and split a binary parenthesization as $\mathcal{A}\otimes\mathcal{B}$.\\
By induction, $\mathcal{A}$ and $\mathcal{B}$ satisfy the appropriate iterated universal properties.\\
Let $f$ be a multilinear map from the original factors to $\mathcal{Y}$.\\
Fixing the factors represented by $\mathcal{B}$ and linearizing the remaining variables produces an element of $\LL(\mathcal{A},\mathcal{Y})$.\\
This assignment is multilinear in the factors represented by $\mathcal{B}$ because the factorizations through $\mathcal{A}$ are unique.\\
Linearizing again produces $L:\mathcal{B}\to\LL(\mathcal{A},\mathcal{Y})$. Set $b_f(\mathbf{a},\mathbf{b})\defeq L(\mathbf{b})(\mathbf{a})$.\\
Then map $b_f$ is bilinear, and induces unique map $\mathcal{A}\otimes\mathcal{B}\to\mathcal{Y}$ by Theorem \ref{thm:tensor-linearizes-bilinear-maps}(ii).\\
Any mediator determines the same two intermediate factorizations and hence the same induced map.\\
Thus the chosen parenthesization is an iterated tensor product.
\item Let $(\mathcal{T},t)$ be an iterated tensor product. Set $\mathcal{S}\defeq\operatorname{span}(t(\mathcal{V}_1\times\cdots\times\mathcal{V}_n))$.\\
The quotient map $\pi:\mathcal{T}\to\mathcal{T}/\mathcal{S}$ is linear by Proposition \ref{prop:quotient-map-properties}(i).\\
Both $\pi$ and the zero map mediate the zero multilinear map because $\pi\circ t=0$.\\
Thus $\pi=0$ by Definition \ref{def:iterated-tensor-product-universal-property}(i). Since $\pi$ is surjective by Proposition \ref{prop:quotient-map-properties}(ii), $\mathcal{T}/\mathcal{S}=0$ and $\mathcal{S}=\mathcal{T}$.
\item Let $(\mathcal{T},t)$ and $(\mathcal{U},u)$ be iterated tensor products. The universal properties produce a unique $F:\mathcal{T}\to\mathcal{U}$ with $F\circ t=u$.\\
They produce a unique $G:\mathcal{U}\to\mathcal{T}$ with $G\circ u=t$.\ Both $GF$ and $I_{\mathcal{T}}$ mediate $t$, while both $FG$ and $I_{\mathcal{U}}$ mediate $u$.\\
Uniqueness implies $GF=I_{\mathcal{T}}$, $FG=I_{\mathcal{U}}$. Any isomorphism preserving decomposable tensors also mediates $u$, so it equals $F$.
\end{enumerate}
\item The zero map is multilinear, and $af+g$ is multilinear when $a\in\mathbb{F}$ and $f,g$ are multilinear.\\
The remaining vector-space axioms hold pointwise because they hold in $\mathcal{W}$. Let $t$ be the canonical multilinear map.\\
Set $\Phi(f)\defeq\widetilde{f}$ by Definition \ref{def:iterated-tensor-product-universal-property}(i). Set $\Psi(T)\defeq T\circ t$. Defining factorisation implies $\Psi\Phi=I$, while uniqueness implies $\Phi\Psi=I$.\\
The maps $\Phi(af+g)$ and $a\Phi(f)+\Phi(g)$ both mediate $af+g$, so uniqueness proves that $\Phi$ is linear. Hence $\Phi$ is an isomorphism.
\item Set $\mathcal{B}\defeq\mathcal{V}_{1:n}^{\otimes}$. Both $\mathcal{A}_r$ and $\mathcal{B}$ are iterated tensor products by (i)(a).\\
Their universal properties produce comparison maps $\alpha_r:\mathcal{A}_r\to\mathcal{B}$ and $\beta_r:\mathcal{B}\to\mathcal{A}_r$.\\
These maps have the displayed actions on decomposable tensors. These maps are inverse isomorphisms by (i)(c).
\item The map $p_\pi:\prod\limits_{i=1}^{n}\mathcal{V}_i\to\mathcal{C}_\pi$ satisfies $p_\pi(\mathbf{v}_1,\ldots,\mathbf{v}_n)\defeq\mathbf{v}_{\pi(1)}\otimes\cdots\otimes\mathbf{v}_{\pi(n)}$ and is multilinear.\\
Its unique mediating map is $\tau_\pi$ by Definition \ref{def:iterated-tensor-product-universal-property}(i). The analogous map for the inverse reordering is inverse to $\tau_\pi$ by (i)(c).
\item The map $m_L:\mathbb{F}\times\mathcal{V}\to\mathcal{V}$ defined by $m_L(r,\mathbf{v})\defeq r\mathbf{v}$ is bilinear. The map $m_R:\mathcal{V}\times\mathbb{F}\to\mathcal{V}$ defined by $m_R(\mathbf{v},r)\defeq r\mathbf{v}$ is bilinear.\\
Hence $\lambda$ and $\rho$ exist by Theorem \ref{thm:tensor-linearizes-bilinear-maps}(ii). Define $u_L:\mathcal{V}\to\mathbb{F}\otimes\mathcal{V}$ by $u_L(\mathbf{v})\defeq1\otimes\mathbf{v}$. Define $u_R:\mathcal{V}\to\mathcal{V}\otimes\mathbb{F}$ by $u_R(\mathbf{v})\defeq\mathbf{v}\otimes1$.\\
Both maps are linear by bilinearity of the canonical tensor maps. Direct computation proves $\lambda u_L=I_{\mathcal{V}}$ and $\rho u_R=I_{\mathcal{V}}$.\\
Maps $u_L\lambda$ and $u_R\rho$ fix decomposable tensors by bilinearity, hence are identity maps by (i)(b). Hence $\lambda$ and $\rho$ are isomorphisms.
\item Choose finite bases $B_i$ of $\mathcal{V}_i$. Set $\mathcal{D}\defeq\{\mathbf{b}_1\otimes\cdots\otimes\mathbf{b}_n:\mathbf{b}_i\in B_i\}$.\\
The set $\mathcal{D}$ is a basis by repeated application of Theorem \ref{thm:tensor-bases-dimension-and-rank}(i). Consequently, $\dim(\mathcal{V}_{1:n}^{\otimes})=\#\mathcal{D}=\prod\limits_{i=1}^{n}\#B_i=\prod\limits_{i=1}^{n}\dim\mathcal{V}_i$.
\end{enumerate}
\hlt{Structural Commentary}: Induction builds products; quotients prove spanning; uniqueness fixes maps.\\
\hlt{Key Insight}: Every parenthesization realizes the same multilinear universal property.
\end{proof}
\end{theorem}

\begin{definition} \label{def:tensor-spaces-type-pq} \hlt{Tensor Spaces of Type $(p,q)$}\\
Let $\dim\mathcal{V}<\infty$.
\begin{enumerate}[label=(\roman*)]
\setlength{\itemsep}{0pt}
\item Set $\mathcal{V}^{\otimes0}\defeq\mathbb{F}$. If $k\in\mathbb{Z}_{>0}$, let $\mathcal{V}^{\otimes k}$ denote the iterated tensor product of $k$ copies of $\mathcal{V}$.
\item Let $(p,q)\in\mathbb{Z}_{\geq0}^{2}\setminus\{(0,0)\}$. Set $\mathcal{T}^{p}_{q}(\mathcal{V})\defeq\mathcal{V}^{\otimes p}\otimes(\mathcal{V}')^{\otimes q}$.\\
Parentheses are suppressed by Theorem \ref{thm:iterated-tensor-products-basic-laws}(iii). Unit factors are identified by part (v).
\item Space $\mathcal{T}^{p}_{q}(\mathcal{V})$ is \hlt{tensor space of type $(p,q)$}. Its $p$ factors from $\mathcal{V}$ are contravariant, $q$ factors from $\mathcal{V}'$ are covariant. $\mathcal{T}^{0}_{0}(\mathcal{V})\defeq\mathbb{F}$.
\item If $p+q\geq1$, set $\mathcal{M}_{p,q}(\mathcal{V})\defeq\operatorname{Mult}_{\mathbb{F}}((\mathcal{V}')^{\times p}\times\mathcal{V}^{\times q};\mathbb{F})$.\\
Its elements are multilinear forms on $p$ arguments from $\mathcal{V}'$ followed by $q$ arguments from $\mathcal{V}$.
\end{enumerate}
\end{definition}

\begin{proposition} \label{prop:tensor-spaces-as-multilinear-functionals} \hlt{Tensor Spaces as Multilinear Functionals}\\
Let $\dim\mathcal{V}=n<\infty$ and $p,q\in\mathbb{Z}_{\geq0}$. Write $\mathbf{v}^{\otimes}\defeq\mathbf{v}_1\otimes\cdots\otimes\mathbf{v}_p$. Write $\boldsymbol{\varphi}^{\otimes}\defeq\varphi_1\otimes\cdots\otimes\varphi_q$.\\
Set $\boldsymbol{\eta}\defeq(\eta_1,\ldots,\eta_p)$ and $\mathbf{x}\defeq(\mathbf{x}_1,\ldots,\mathbf{x}_q)$. In an argument list, $(\boldsymbol{\eta},\mathbf{x})$ denotes the concatenation of these tuples.\\
Set $P_{\mathbf{v}}(\boldsymbol{\eta})\defeq\prod\limits_{a=1}^{p}\eta_a(\mathbf{v}_a)$ and $Q_{\boldsymbol{\varphi}}(\mathbf{x})\defeq\prod\limits_{b=1}^{q}\varphi_b(\mathbf{x}_b)$. An empty tensor or product below equals $1$.
\begin{enumerate}[label=(\roman*)]
\setlength{\itemsep}{0pt}
\item $\dim\mathcal{T}^{p}_{q}(\mathcal{V})=n^{p+q}$.
\item If $p+q\geq1$, a canonical isomorphism $\Theta_{p,q}:\mathcal{T}^{p}_{q}(\mathcal{V})\xrightarrow{\sim}\mathcal{M}_{p,q}(\mathcal{V})$ is determined by $\Theta_{p,q}(\mathbf{v}^{\otimes}\otimes\boldsymbol{\varphi}^{\otimes})(\boldsymbol{\eta},\mathbf{x})=P_{\mathbf{v}}(\boldsymbol{\eta})Q_{\boldsymbol{\varphi}}(\mathbf{x})$.
\item Suppose $p,q\geq1$ and $(i,j)\in\{1,\ldots,p\}\times\{1,\ldots,q\}$. A subscript $\widehat{i}$ or $\widehat{j}$ denotes omission from the corresponding tensor.\\
Set $\mathbf{z}_{ij}\defeq\mathbf{v}_{\widehat{i}}^{\otimes}\otimes\boldsymbol{\varphi}_{\widehat{j}}^{\otimes}$. A unique linear map $C_{ij}:\mathcal{T}^{p}_{q}(\mathcal{V})\to\mathcal{T}^{p-1}_{q-1}(\mathcal{V})$ satisfies $C_{ij}(\mathbf{v}^{\otimes}\otimes\boldsymbol{\varphi}^{\otimes})=\varphi_j(\mathbf{v}_i)\mathbf{z}_{ij}$.
\item If $n\geq1$, an isomorphism $R:\mathcal{V}\otimes\mathcal{V}'\xrightarrow{\sim}\LL(\mathcal{V})$ satisfies $R(\mathbf{v}\otimes\varphi)(\mathbf{x})\defeq\varphi(\mathbf{x})\mathbf{v}$. On $\mathcal{T}^{1}_{1}(\mathcal{V})$, $C_{11}=\operatorname{tr}\circ R$.
\end{enumerate}
\begin{proof}
\begin{enumerate}[label=(\roman*)]
\setlength{\itemsep}{0pt}
\item The equality $\dim\mathcal{V}'=n$ holds by Theorem \ref{thm:dual-basis-and-dimension}(ii). If $p=q=0$, then $\mathcal{T}^{0}_{0}(\mathcal{V})=\mathbb{F}$ by Definition \ref{def:tensor-spaces-type-pq}(iv).\\
Hence $\dim\mathcal{T}^{0}_{0}(\mathcal{V})=1=n^0$, where $n^0\defeq1$. If $p+q\geq1$, then $\mathcal{T}^{p}_{q}(\mathcal{V})$ has $p$ factors of dimension $n$ and $q$ factors of dimension $n$.\\
Its dimension is $n^{p+q}$ by Theorem \ref{thm:iterated-tensor-products-basic-laws}(vi) and Definition \ref{def:tensor-spaces-type-pq}(ii).
\item Target of $\Theta_{p,q}$ is vector space by Theorem \ref{thm:iterated-tensor-products-basic-laws}(ii). Displayed formula defines multilinear form in $\eta_1,\ldots,\eta_p,\mathbf{x}_1,\ldots,\mathbf{x}_q$.\\
Varying $\mathbf{v}_1,\ldots,\mathbf{v}_p,\varphi_1,\ldots,\varphi_q$ defines a multilinear assignment into that target.\\
Unique linear map $\Theta_{p,q}$ therefore exists by Theorem \ref{thm:iterated-tensor-products-basic-laws}(ii). Set $\mathcal{X}\defeq(\mathcal{V}')^{\otimes p}\otimes\mathcal{V}^{\otimes q}$.\\
The codomain is canonically isomorphic to $\mathcal{X}'$ by Theorem \ref{thm:iterated-tensor-products-basic-laws}(ii). The domain has dimension $n^{p+q}$ by part (i).\\
The equality $\dim\mathcal{X}=n^{p+q}$ holds by Theorem \ref{thm:iterated-tensor-products-basic-laws}(vi). The codomain has dimension $\dim\mathcal{X}'=\dim\mathcal{X}$ by Theorem \ref{thm:dual-basis-and-dimension}(ii).\\
If $n=0$, this dimension is zero because $p+q\geq1$, so $\Theta_{p,q}$ is an isomorphism.\\
Suppose $n\geq1$. Choose a basis $(\mathbf{e}_1,\ldots,\mathbf{e}_n)$ and its dual basis $(\varepsilon_1,\ldots,\varepsilon_n)$.\\
For $I=(i_1,\ldots,i_p)$, set $\mathbf{e}_I^{\otimes}\defeq\mathbf{e}_{i_1}\otimes\cdots\otimes\mathbf{e}_{i_p}$. For $J=(j_1,\ldots,j_q)$, set $\varepsilon_J^{\otimes}\defeq\varepsilon_{j_1}\otimes\cdots\otimes\varepsilon_{j_q}$. Set $\mathbf{b}_{I,J}\defeq\mathbf{e}_I^{\otimes}\otimes\varepsilon_J^{\otimes}$.\\
The elements $\mathbf{b}_{I,J}$ form a basis by repeated application of Theorem \ref{thm:tensor-bases-dimension-and-rank}(i).\\
For $K=(k_1,\ldots,k_p)$, set $\delta_{K,I}\defeq\prod\limits_{a=1}^{p}\delta_{k_a,i_a}$. For $L=(\ell_1,\ldots,\ell_q)$, set $\delta_{J,L}\defeq\prod\limits_{b=1}^{q}\delta_{j_b,\ell_b}$.\\
Set $\boldsymbol{\varepsilon}_K\defeq(\varepsilon_{k_1},\ldots,\varepsilon_{k_p})$. Set $\mathbf{e}_L\defeq(\mathbf{e}_{\ell_1},\ldots,\mathbf{e}_{\ell_q})$. The defining formula implies $\Theta_{p,q}(\mathbf{b}_{I,J})(\boldsymbol{\varepsilon}_K,\mathbf{e}_L)=\delta_{K,I}\delta_{J,L}$.\\
If $z\in\ker\Theta_{p,q}$, evaluation at $(\boldsymbol{\varepsilon}_I,\mathbf{e}_J)$ equals both zero and the coefficient of $\mathbf{b}_{I,J}$ in $z$.\\
Thus every coefficient is zero, so $\Theta_{p,q}$ is injective. It is an isomorphism by Corollary \ref{cor:equal-dimension-kernel-range}(ii).
\item Let $\Gamma_{ij}$ be the rule in part (iii). In the selected variables $\mathbf{v}_i$ and $\varphi_j$, linearity follows from bilinearity of evaluation.\\
In every remaining variable, linearity holds by Theorem \ref{thm:iterated-tensor-products-basic-laws}(i)(a).\\
Thus $\Gamma_{ij}$ is multilinear, and a unique linear map $C_{ij}$ with the required action exists by part (ii) of that theorem.
\item The rule $(\mathbf{v},\varphi)\mapsto(\mathbf{x}\mapsto\varphi(\mathbf{x})\mathbf{v})$ is bilinear. Thus $R$ exists by Theorem \ref{thm:tensor-linearizes-bilinear-maps}(ii).\\
Choose a basis $(\mathbf{e}_1,\ldots,\mathbf{e}_n)$ and its dual basis $(\varepsilon_1,\ldots,\varepsilon_n)$. The map $R$ sends $\mathbf{e}_a\otimes\varepsilon_b$ to the operator whose matrix is $E_{a,b}$.\\
The source elements form a basis by Theorem \ref{thm:tensor-bases-dimension-and-rank}(i). The matrices $E_{a,b}$ form a basis by Proposition \ref{prop:matrix-spaces-are-vector-spaces}(ii).\\
Matrix representation identifies this basis with an operator basis by Theorem \ref{thm:linear-maps-and-matrices-same-data}(i).\\
Hence $R$ is an isomorphism. Write $\mathbf{v}=\sum\limits_{a=1}^{n}x_a\mathbf{e}_a$. The diagonal entries of $[R(\mathbf{v}\otimes\varphi)]$ are $x_a\varphi(\mathbf{e}_a)$.\\
The identity $\operatorname{tr}(R(\mathbf{v}\otimes\varphi))=\sum\limits_{a=1}^{n}x_a\varphi(\mathbf{e}_a)$ holds by Definition \ref{def:operator-trace-and-determinant}(i).\\
Sum equals $\varphi(\mathbf{v})=C_{11}(\mathbf{v}\otimes\varphi)$ by linearity and part (iii). Maps $\operatorname{tr}\circ R$ and $C_{11}$ are linear by Proposition \ref{prop:trace-arithmetic} and part (iii).\\
They agree on the spanning decomposable tensors by Theorem \ref{thm:iterated-tensor-products-basic-laws}(i)(b), so they are equal.
\end{enumerate}
\hlt{Structural Commentary}: Bases establish dimension and duality; evaluation defines contractions.\\
\hlt{Key Insight}: Contraction evaluates one vector--covector pair and preserves the remaining slots.
\end{proof}
\end{proposition}

\begin{definition} \label{def:contraction-map} \hlt{Contraction}\\
Let $\dim\mathcal{V}<\infty$ and let $p,q\in\mathbb{Z}_{>0}$. The maps $C_{ij}$ in Proposition \ref{prop:tensor-spaces-as-multilinear-functionals}(iii) are the \hlt{contractions} of $\mathcal{T}^{p}_{q}(\mathcal{V})$.\\
The \hlt{standard contraction} is $C\defeq C_{11}$.
\end{definition}

\begin{example} \label{ex:rank-one-operator-contraction-trace} \hlt{Rank-One Operators and Trace}\\
Let $n\defeq\dim\mathcal{V}\in\mathbb{Z}_{>0}$ and let $(\mathbf{e}_1,\ldots,\mathbf{e}_n)$ be a basis. Let $(\varepsilon_1,\ldots,\varepsilon_n)$ be the dual basis.\\
By Proposition \ref{prop:tensor-spaces-as-multilinear-functionals}(iv), $R(\mathbf{e}_a\otimes\varepsilon_b)$ has matrix $E_{a,b}$ for $a,b\in\{1,\ldots,n\}$. Its contraction records the trace: $C(\mathbf{e}_a\otimes\varepsilon_b)=\delta_{a,b}=\operatorname{tr}(E_{a,b})$.
\end{example}
